\subsection{测试小节一}

本节测试第一轮未覆盖的定理环境.

\begin{axiom}{测试公理}{ax:test}
    这是测试用的公理.
\end{axiom}

\begin{hypothesis}{测试假设}{hyp:test}
    这是测试用的假设.
\end{hypothesis}

\begin{metatheorem}{测试元定理}{mt:test}
    这是测试用的元定理.
\end{metatheorem}

\begin{criteria}{测试准则}{cri:test}
    这是测试用的准则.
\end{criteria}

\begin{problem}{测试问题}{prob:test}
    这是测试用的问题.
\end{problem}

\begin{exercise}{测试习题}{exr:test}
    这是测试用的习题, 验证独立编号与品红配色.
\end{exercise}

\begin{solution}
    这是测试用的解答, 紧跟在习题之后, 验证题后即答布局. 本解答引用讲义内的定理 \noteref{thm:test}.
\end{solution}

\begin{exercise}{另一道习题}{exr:test2}
    这是第二道习题, 检验编号是否连续递增.
\end{exercise}

\begin{algorithm}{测试算法}{alg:test}
    输入数据, 执行若干步骤, 然后返回结果.
\end{algorithm}

\begin{convention}{测试约定}{conv:test}
    这是测试用的约定.
\end{convention}

\begin{alarm}{测试警示}{alr:test}
    这是测试用的警示.
\end{alarm}

\begin{sketch}{测试分析}{sk:test}
    这是测试用的分析.
\end{sketch}

\begin{answer}
    这是测试用的回答.
\end{answer}

\begin{hint}[测试提示]
    这是测试用的提示内容.
\end{hint}

\begin{code}[Python 测试]
def square(x):
    return x * x
\end{code}

环境引用测试: \autoref{ax:test}, \autoref{hyp:test}, \autoref{mt:test}, \autoref{cri:test}, \autoref{prob:test}, \autoref{exr:test}, \autoref{alg:test}, \autoref{conv:test}, \autoref{alr:test}, \autoref{sk:test}.
